\chapter{Statikus assetek és frontend-integráció}

\noindent\textbf{Tanulási út: Gyors út: böngészős ág.} Most olvasd, ha static assetet vagy progresszív böngészős enhancementet adsz hozzá.\par\medskip

A webalkalmazásnak nem csak dinamikus HTML- és JSON-válaszai vannak. CSS, ikon, kép, font, source map vagy egy külön frontend build artifactjai más életciklusú fájlok: a release részeként érkeznek, tartalmuk ideális esetben immutable, és a böngésző agresszíven cache-elheti őket.

A Velran ezért a statikus asseteket külön, \textbf{read-only static rootból} szolgálja ki. Ez nem ugyanaz a terület, mint az AppFs írható data rootja, és nem ugyanaz a fogalom, mint a felhasználó által feltöltött média.

\noindent\textbf{Első olvasás:} tanuld meg a három fájléletciklust, a fingerprintelt cache-elést és a V1 JavaScript/CSP határt. A szerver-vs-proxy kiszolgálási példák deployment referenciaként maradhatnak későbbre.

\section{Három külön fájléletciklus}

Egy production alkalmazásban legalább három fájlcsoportot érdemes külön kezelni:

\begin{center}
\begin{tabularx}{\textwidth}{@{}lYYY@{}}
\toprule
 & Példa & Írhatóság & Életciklus \\
\midrule
release & CSS, font, logo & read-only & deploykor cserélődik \\
runtime data & export, generált fájl & AppFs policy szerint & alkalmazás futása közben \\
user media & feltöltött kép & kontrolláltan írható & üzleti adattal együtt \\
\bottomrule
\end{tabularx}
\end{center}

Tipikus könyvtárszerkezet:

\begin{lstlisting}
/usr/local/bin/velran-server
/usr/local/bin/velran-cli
/srv/myapp/releases/2026-09-04/
  main.vrn
  migrations/
  public/
    css/
    img/
    fonts/
/srv/myapp/data/
/run/secrets/velran/
\end{lstlisting}

A \code{public/} legyen deployment artifact. A \code{data/} ezzel szemben runtime állapot. Ha a kettőt összekevered, egy upload felülírhat release-fájlt, vagy egy rollback véletlenül felhasználói adatot cserélhet vissza.

\section{A static root konfigurálása}

A server configban a statikus namespace külön szekció:

\begin{lstlisting}
[static_assets]
root = "/srv/myapp/current/public"
url_prefix = "/assets/"
max_asset_bytes = 8388608
max_age_secs = 300
immutable_max_age_secs = 31536000
precompressed = true
\end{lstlisting}

Ugyanez CLI-ből is beállítható. A fontosabb opciók:

\begin{lstlisting}
--static-root
--static-url-prefix
--max-static-asset-bytes
--static-max-age-secs
--static-immutable-max-age-secs
\end{lstlisting}

A static prefix rezervált route namespace. Ha az URL prefix \code{/assets/}, az application ne deklaráljon ezzel ütköző dinamikus vagy statikus route-ot. Az ütközést a szerver fail-closed módon kezelje: jobb startupkor hibázni, mint requestenként bizonytalan route-választást végezni.

\section{Fingerprint: a hosszú cache kulcsa}

A statikus asset tartalma és URL-je legyen összekötve. Például:

\begin{lstlisting}
public/
  css/app.01234567.css
  img/logo.a91f028c.svg
  fonts/ui-f1e2d3c4.woff2
\end{lstlisting}

A Velran fingerprinteltnek tekinti azt a fájlnevet, amelyben legalább 8 hexadecimális karakterből álló, ponttal, kötőjellel vagy underscore-ral határolt token szerepel. Ilyen fájlnál a response hosszú immutable cache-t kaphat:

\begin{lstlisting}
Cache-Control: public, max-age=31536000, immutable
\end{lstlisting}

Nem fingerprintelt fájlnál a biztonságosabb alap rövid és revalidálható:

\begin{lstlisting}
Cache-Control: public, max-age=300, must-revalidate
\end{lstlisting}

A lényeg: \textbf{ne ugyanazon URL mögött cseréld le a tartalmat egyéves cache mellett}. Új CSS buildhez új fingerprintelt fájlnév tartozzon. A régi fájl akár maradhat is a release-ben addig, amíg a régi HTML-reprezentációk kifuthatnak.

\section{Egy minimális asset build pipeline}

A Velran nem akar frontend build rendszert előírni. A pipeline feladata egyszerűen az, hogy determinisztikus outputot készítsen a \code{public/} könyvtárba.

Egy kis alkalmazásnál a forrás lehet például:

\begin{lstlisting}
assets-src/
  app.css
  logo.svg
public/
\end{lstlisting}

A release build lépései:

\begin{enumerate}
  \item tiszta \code{public/} könyvtár létrehozása;
  \item CSS/font/kép build vagy másolás;
  \item content hash képzése;
  \item fingerprintelt fájlnév létrehozása;
  \item opcionális Brotli/gzip siblingek készítése;
  \item az alkalmazásban a fingerprintelt URL használata.
\end{enumerate}

A hash nem security token. A feladata az, hogy a fájlnév a tartalom verzióját azonosítsa.

\section{ETag és conditional GET}

Minden statikus reprezentáció saját \code{ETag} értéket kap. A böngésző később küldhet:

\begin{lstlisting}
If-None-Match: "velran-..."
\end{lstlisting}

Ha a reprezentáció nem változott, a szerver \code{304 Not Modified} választ adhat body nélkül. GET mellett HEAD is támogatott.

A static request nem hoz létre sessiont és nem küld session cookie-t. Ez fontos teljesítmény- és security tulajdonság: egy CSS vagy font letöltése ne váljon session-lifecycle eseménnyé.

\section{Precompressed Brotli és gzip}

A Velran server request közben nem tömöríti a statikus fájlt. A release pipeline opcionálisan előre elkészítheti:

\begin{lstlisting}
app.01234567.css
app.01234567.css.br
app.01234567.css.gz
\end{lstlisting}

A kliens \code{Accept-Encoding} értéke alapján a szerver a Brotli, majd gzip, végül az eredeti reprezentáció közül választ. Tömörített válasznál a header például:

\begin{lstlisting}
Content-Encoding: br
Vary: Accept-Encoding
\end{lstlisting}

Ennek két előnye van: nincs requestenkénti compression CPU-költség, és a web request nem nyit dinamikus tömörítési erőforrás-felületet.

\section{MIME, nosniff és path confinement}

A server ismert webes kiterjesztésekhez MIME típust rendel, ismeretlen extensionnél pedig \code{application/octet-stream} típust használ. A válasz mindig kap:

\begin{lstlisting}
X-Content-Type-Options: nosniff
\end{lstlisting}

A static path szándékosan szigorú. Jó:

\begin{lstlisting}
/assets/css/app.01234567.css
\end{lstlisting}

Tiltott többek között a \code{..}, encoded traversal, backslash, üres path komponens és control karakter. A V1 static pathban percent-encoding sincs; a build pipeline használjon URL-safe fájlneveket.

Linuxon a static root ugyanahhoz a confinement szemlélethez igazodik, mint az AppFs: a kiszolgálás nem válhat tetszőleges host filesystem olvasássá symlink vagy path traversal segítségével.

\section{CSS, kép és font a Velran HTML-ben}

A V1 CSP same-origin style, image, font és media asseteket enged. Egy release-ben ezért természetes minta:

\begin{lstlisting}
/assets/css/app.01234567.css
/assets/img/logo.a91f028c.svg
/assets/fonts/ui-f1e2d3c4.woff2
\end{lstlisting}

A HTML-ben mindig a release-hez tartozó fingerprintelt URL jelenjen meg. A sablon ne felhasználói inputból állítson elő asset URL-t; az assetnév a build része, nem request adat.

\section{Fontos V1-határ: JavaScript kiszolgálható, de Velran HTML-ben nem fut}

A static server ismeri a \code{.js} és \code{.mjs} MIME típust, ezért ilyen build artifact fizikailag kiszolgálható. Ez azonban \textbf{nem jelenti azt, hogy a Velran által renderelt HTML scriptet futtathat}.

A V1 HTML security modellben:

\begin{itemize}
  \item a CSP nem nyit \code{script-src} engedélyt;
  \item a compiler a \code{<script>} markupot elutasítja;
  \item nincs inline-script vagy ``unsafe'' escape hatch.
\end{itemize}

Ez tudatos security határ. A könyv előző fejezetének JavaScript \code{fetch} példái ezért kétféleképpen értelmezendők:

\begin{enumerate}
  \item külön frontend alkalmazás fut a saját originjén, és Velran JSON API-t hív;
  \item a kód a browser API kliens viselkedését szemlélteti, nem azt, hogy V1 Velran HTML-be script tag írható.
\end{enumerate}

Ha az alkalmazás klasszikus Velran szerver-renderelt oldal, építs HTML + CSS + typed form/action modellre. Ha teljes JavaScript frontend szükséges, kezeld azt külön frontendként, és a Velran legyen typed JSON backend.

\section{Velran szolgálja ki vagy Apache/Nginx?}

Két működő deployment modell van.

\subsection*{1. Velran-owned static namespace}

Ez legyen az alapértelmezett. A reverse proxy minden requestet a Velran szervernek ad, és a Velran kezeli az \code{/assets/} namespace-t is:

\begin{lstlisting}
Internet
   |
Apache / Nginx :443
   |
127.0.0.1:8080 Velran
   |-- application routes
   `-- /assets/*
\end{lstlisting}

Előnye, hogy a fingerprint-, ETag-, precompressed-, nosniff- és path-policy egy helyen marad. Kis és közepes alkalmazásnál ez általában bőven elég.

\subsection*{2. A reverse proxy szolgálja ki közvetlenül}

Nagy forgalomnál dönthetsz úgy, hogy az Apache vagy Nginx közvetlenül a read-only release \code{public/} könyvtárból szolgál. Ekkor viszont \textbf{neked kell reprodukálnod} a fontos contractot:

\begin{itemize}
  \item read-only, külön static root;
  \item traversal/symlink elleni biztonságos path-kezelés;
  \item helyes MIME + \code{nosniff};
  \item fingerprintelt fájlokra hosszú immutable cache;
  \item nem fingerprintelt fájlokra rövid/revalidálható cache;
  \item Brotli/gzip esetén helyes \code{Vary};
  \item a namespace ne ütközzön application route-tal.
\end{itemize}

Ez már operátori optimalizáció, nem szükséges ahhoz, hogy az alkalmazás helyes legyen.

\section{Nginx minta közvetlen static kiszolgálásra}

Egy egyszerű kiindulás:

\Needspace{15\baselineskip}
\begin{lstlisting}
server {
    listen 443 ssl;
    server_name app.example;

    location /assets/ {
        alias /srv/myapp/current/public/;
        try_files $uri =404;
        add_header X-Content-Type-Options nosniff always;
    }

    location / {
        proxy_pass http://127.0.0.1:8080;
        proxy_set_header Host $host;
        proxy_set_header X-Forwarded-Proto https;
    }
}
\end{lstlisting}

Ez csak topológiai alapminta. A production konfigurációban a cache policy-t, precompressed fájlokat és symlink/release szabályokat is ugyanahhoz az immutable deployment modellhez kell igazítani.

\section{Apache minta közvetlen static kiszolgálásra}

Apache esetén a gondolat ugyanaz:

\Needspace{15\baselineskip}
\begin{lstlisting}
Alias /assets/ /srv/myapp/current/public/

<Directory /srv/myapp/current/public/>
    Require all granted
    Options -Indexes
    AllowOverride None
    Header always set X-Content-Type-Options "nosniff"
</Directory>

ProxyPass        / http://127.0.0.1:8080/
ProxyPassReverse / http://127.0.0.1:8080/
\end{lstlisting}

Ha Apache közvetlenül szolgálja az asseteket, az \code{Alias} szabálynak a proxy catch-all előtt kell érvényesülnie, és a cache headereket is explicit konfigurálni kell.

\section{Deploy: asset és alkalmazáskód ugyanahhoz a release-hez tartozzon}

A legkönnyebben elkerülhető assethiba az, amikor az új HTML már új CSS fájlnevet hivatkozik, de a static rootban még a régi release van -- vagy fordítva.

Ezért a helyes sorrend:

\begin{lstlisting}
build application
-> build/fingerprint assets
-> optional .br/.gz
-> immutable release directory
-> preflight
-> atomic current switch
-> controlled restart/reload
\end{lstlisting}

Az alkalmazáskód és a \code{public/} könyvtár ugyanazon release artifact része legyen. Így rollbackkor mindkettő együtt áll vissza.

\section{Ellenőrzés curl-lel}

Deploy után legalább ezeket érdemes ellenőrizni:

\begin{lstlisting}
curl -I https://app.example/assets/css/app.01234567.css
curl -I -H 'Accept-Encoding: br' \
  https://app.example/assets/css/app.01234567.css
curl -I -H 'If-None-Match: "velran-..."' \
  https://app.example/assets/css/app.01234567.css
\end{lstlisting}

Ellenőrizd a \code{Content-Type}, \code{Cache-Control}, \code{ETag}, \code{X-Content-Type-Options}, valamint tömörített reprezentációnál a \code{Content-Encoding} és \code{Vary} headert.

\section{Gyakorlati döntési szabály}

\begin{center}
\begin{tabularx}{\textwidth}{@{}lY@{}}
\toprule
Helyzet & Ajánlott megoldás \\
\midrule
szerver-renderelt Velran oldal & Velran static root + CSS/image/font \\
kis/közepes production & reverse proxy mindent Velrannak továbbít \\
nagy statikus forgalom & Apache/Nginx direct static, azonos policy-vel \\
teljes JavaScript frontend & külön frontend deployment + Velran JSON API \\
user upload & AppFs/media pipeline, nem static root \\
\bottomrule
\end{tabularx}
\end{center}

A legfontosabb szabály: \textbf{a statikus asset deployment artifact, nem alkalmazásadat}. Ha ezt a határt megtartod, a cache, rollback, backup és security modell egyszerre lesz egyszerűbb.

\section{Ellenőrzött companion példa}
Az aktuális static-asset companion az \path{examples/static-assets/app.vrn}. A release asseteket tartsd read-only módon, az AppFs/runtime állapottól elkülönítve, a fejezet deployment modelljének megfelelően.

\section{Gyakorlat: osztályozd a frontend minden fájlját}
\paragraph{Gyakorlási mód: támogatott.} Használd az előszóban meghatározott támogatott módszert.

Egy page által használt minden fájlt sorolj immutable deployment asset, generált HTML, JSON response vagy módosuló alkalmazásadat kategóriába. Ellenőrizd, hogy a static root alatt kizárólag immutable deployment asset található.

\section{Tipikus hiba: a frontend-kényelem eltünteti a trust boundaryt}
Egy JavaScript frontend nem változtatja meg a szerveroldali autorizáció követelményeit. A kliensoldali validáció és elrejtett vezérlők UX-eszközök; minden authoritative validáció, autorizáció és módosítási szabály a szerver határán marad.

\section{Folyamatos mintaprojekt: Secure Notes}
Adj kis stylesheetet és szükség esetén progresszív JavaScript enhancementet a jegyzet UI-hoz. Az alkalmazás maradjon helyes kliensoldali ellenőrzésekbe vetett bizalom nélkül is, a felhasználói tartalom pedig maradjon a static asset fán kívül.
